\chapter{Statistiques et lecture critique}
\label{exercices:10}
\section*{Exercice 10.1 -- Tous les indicateurs}
Les notes et effectifs sont :
\[
\begin{array}{c|rrrrr}\text{Note}&8&10&12&14&16\\\hline\text{Effectif}&2&3&6&4&1\end{array}
\]
Calculer l'effectif total, la fréquence des notes au moins égales à $14$, la moyenne, la médiane et l'étendue.

\section*{Exercice 10.2 -- Choisir et construire un graphique}
Des durées sont regroupées dans $[0;10[$, $[10;20[$, $[20;30[$ avec effectifs $5,9,6$. Construire l'histogramme. Calculer l'angle du secteur de la deuxième classe dans un diagramme circulaire. Peut-on retrouver exactement la moyenne ?

\section*{Exercice 10.3 -- Moyenner des moyennes}
Un groupe de $12$ élèves a une moyenne de $10$ et un groupe de $18$ élèves une moyenne de $15$. Calculer la moyenne commune. Expliquer l'erreur de la réponse $12{,}5$.

\section*{Exercice 10.4 -- Valeur extrême et lecture critique}
Comparer moyenne, médiane et étendue de $2,3,4,5,6$ et de $2,3,4,5,26$. Quel indicateur change le moins ? Un graphique dont l'axe commence à $20$ peut-il être trompeur pour comparer des valeurs $21$ et $22$ ?
